\documentclass[11pt,a4paper]{article}

\usepackage[margin=2.5cm]{geometry}
\usepackage{amssymb}
\usepackage{booktabs}
\usepackage{tabularx}
\usepackage{multirow}
\usepackage{hyperref}
\usepackage[table]{xcolor}
\usepackage{enumitem}
\usepackage{lmodern}
\usepackage[T1]{fontenc}

\hypersetup{
  colorlinks=true,
  linkcolor=blue!60!black,
  urlcolor=blue!60!black,
}

\setlength{\parindent}{0pt}
\setlength{\parskip}{0.6em}

\title{Comparing the RISC-V ISA Profiles of\\
  Two XMSS C~Implementations}
\author{Wrenna Robson\thanks{With Claude (Anthropic), which co-authored
  the analysis tooling and report text.}}
\date{April 2026}

\begin{document}
\maketitle

\section{Context}

Two independent C~implementations of XMSS (RFC~8391) have been
profiled on RISC-V~64 in companion reports:

\begin{itemize}[nosep]
  \item \emph{RISC-V ISA Profile of the XMSS C Implementation} --- our
    from-scratch C~implementation in \texttt{impl/c/}, designed to be
    structurally portable to Jasmin (no VLAs, no heap allocation, no
    function pointers).
  \item \emph{RISC-V ISA Profile of the Upstream XMSS Reference
    Implementation} --- the authoritative reference at
    \href{https://github.com/XMSS/xmss-reference}{\texttt{XMSS/xmss-reference}}.
\end{itemize}

This report compares the two profiles side by side.  Both were produced
with the same toolchain (\texttt{riscv64-linux-gnu-gcc}~13.3.0,
\texttt{-march=rv64gc -mabi=lp64d -O3}), the same classifier (a
mnemonic lookup table generated from the \texttt{riscv-opcodes}
database), and the same per-object-file and per-extension tabulation.
Any differences between the two profiles are therefore attributable
to differences in \emph{source code}, not to toolchain or methodology.

\section{Archives under comparison}

\subsection{What is \emph{in} each archive}

\begin{tabular}{@{}llc@{}}
\toprule
Component                    & \texttt{impl/c/libxmss.a} & \texttt{libxmss-ref.a} \\
\midrule
SHA-256 / SHA-512 core       & \checkmark (bundled)      & --- (OpenSSL \texttt{libcrypto}) \\
SHAKE128 / SHAKE256          & \checkmark (bundled)      & \checkmark (bundled \texttt{fips202.c}) \\
WOTS+                        & \checkmark                & \checkmark \\
L-tree                       & \checkmark                & \checkmark (inside \texttt{xmss\_commons}) \\
Treehash / BDS state         & \checkmark                & \checkmark (inside \texttt{xmss\_core\_fast}) \\
BDS state (de)serialisation  & \checkmark                & --- (in-memory only) \\
XMSS / XMSS-MT core          & \checkmark                & \checkmark \\
Top-level API                & \checkmark                & \checkmark \\
Parameter derivation         & \checkmark                & \checkmark \\
\bottomrule
\end{tabular}

\subsection{Structural asymmetry: SHA-2}
\label{sec:sha2-asym}

The decisive asymmetry is SHA-2.  The reference's \texttt{hash.c}
calls \texttt{SHA256}/\texttt{SHA512} from OpenSSL; its SHA-2
compression code lives in \texttt{libcrypto} and is \emph{not} present
in the archive.  Our implementation ships its own SHA-2 in pure C
(\texttt{sha2\_local.c}).

As a result, direct counts are not apples-to-apples.  A like-for-like
comparison is only possible by netting out our SHA-2 layer (or by
measuring only the algorithm layer).  Both views are shown below.

\subsection{Module mapping}

The two code bases carve the XMSS algorithm differently.
Table~\ref{tab:mapping} establishes the correspondence used in the
per-module comparisons.

\begin{table}[ht]
\centering
\caption{Module mapping between the two implementations.}
\label{tab:mapping}
\begin{tabularx}{\linewidth}{@{}XXX@{}}
\toprule
Role & \texttt{impl/c/} & Reference \\
\midrule
SHA-2 core          & \texttt{sha2\_local.c.o}       & (in libcrypto) \\
Keccak / SHAKE      & \texttt{shake\_local.c.o}      & \texttt{fips202.o} \\
Hash wrappers       & \texttt{xmss\_hash.c.o}        & \texttt{hash.o} \\
ADRS manipulation   & \texttt{address.c.o}           & \texttt{hash\_address.o} \\
WOTS+               & \texttt{wots.c.o}              & \texttt{wots.o} \\
L-tree              & \texttt{ltree.c.o}             & \multirow{3}{\linewidth}{\texttt{xmss\_commons.o}} \\
Treehash            & \texttt{treehash.c.o}          & \\
BDS state (runtime) & \texttt{bds.c.o}               & \\
BDS serialisation   & \texttt{bds\_serialize.c.o}    & (not present) \\
XMSS / XMSS-MT core & \texttt{xmss.c.o}, \texttt{xmss\_mt.c.o} & \texttt{xmss\_core\_fast.o} \\
Top-level API       & (inside \texttt{xmss.c.o})     & \texttt{xmss.o} \\
Parameters          & \texttt{params.c.o}            & \texttt{params.o} \\
Utilities           & \texttt{utils.c.o}             & \texttt{utils.o} \\
\bottomrule
\end{tabularx}
\end{table}

\section{Headline numbers}

Table~\ref{tab:headline} gives the top-level counts.

\begin{table}[ht]
\centering
\caption{Headline comparison with \texttt{-march=rv64gc}.}
\label{tab:headline}
\begin{tabular}{@{}lrrr@{}}
\toprule
Metric                        & \texttt{impl/c/}      & Reference     & Reference $-$ impl/c \\
\midrule
Object files                  & 13                    & 9             & $-$4 \\
Total instructions            & 9562                  & 7620          & $-$1942 \\
\quad I                       & 9505 \ (99.40\%)      & 7541 \ (98.96\%) & $-$1964 \\
\quad M                       &   57 \ (0.60\%)       &   79 \ (1.04\%)  & $+$22 \\
\quad Zbb (absent in rv64gc)  &    0                  &    0          & --- \\
C-encoded ratio               &   48\%                &   53\%         & $+$5 pp \\
Unique mnemonics              &   52                  &   54           & $+$2 \\
\bottomrule
\end{tabular}
\end{table}

On absolute count, the reference archive is 1942 instructions smaller.
Almost all of that gap, 1953 instructions, is our
\texttt{sha2\_local.c.o}: once SHA-2 is netted out, the two archives
are within 11 instructions of each other (7609 vs 7620).  In other
words, the archives are the same size up to the SHA-2 layer that the
reference excludes by construction.

The extension surface is also nearly identical: both use only I~+~M,
with A, F, D, Zba, Zbb, Zbs, Zicsr, and Zifencei absent from both.

\section{Where the M instructions are}

Both implementations fall well below 2\% M usage, but the composition
differs.

\begin{table}[ht]
\centering
\caption{M-extension mnemonic breakdown.}
\label{tab:m-cmp}
\begin{tabular}{@{}lrrll@{}}
\toprule
Mnemonic        & \texttt{impl/c/} & Reference & Compiler emits it for                              \\
\midrule
\texttt{mulw}   & 37               & 68        & 32-bit index $\times$ parameter \\
\texttt{mul}    & 17               & 6         & 64-bit offset / pointer arithmetic \\
\texttt{divuw}  &  3               & 1         & parameter derivation ($h/d$) \\
\texttt{divu}   &  0               & 2         & SHAKE absorb/squeeze offset (reference only) \\
\texttt{remu}   &  0               & 2         & SHAKE block-length modulus (reference only) \\
\midrule
Total           & 57               & 79        & \\
\bottomrule
\end{tabular}
\end{table}

The reference's 22 extra M-instructions come from two differences:

\begin{itemize}[nosep]
  \item \textbf{SHAKE implementation style.}  The reference's
    \texttt{fips202.c} uses ordinary \texttt{/}~and~\texttt{\%}
    arithmetic on byte counts in the absorb/squeeze loops
    (\texttt{inlen \% rate}, \texttt{nblocks = outlen/rate}).  Our
    SHAKE implementation works in fixed block counts with
    power-of-two-friendly arithmetic, so the compiler never emits
    \texttt{divu}/\texttt{remu}.
  \item \textbf{Index-multiplication density.}  The reference uses
    \texttt{mulw} 68~times; we use it 37~times.  The gap mostly
    lives in \texttt{xmss\_core\_fast.o} (59~M uses in the reference)
    versus the combined \texttt{xmss\_mt.o}\,+\,\texttt{xmss.o}\,+\,\texttt{bds.o}
    (21~M uses in our code).  The reference computes signature-segment
    offsets (\texttt{i * params->n}, \texttt{i * 2*n}) inline in many
    places; our code factors common offsets into locals and carries
    them through the loops, so the compiler re-uses a single
    register-materialised offset instead of re-multiplying.
\end{itemize}

Neither pattern indicates any structural dependence on M: an
RV64I-only target would compile the same C, with shift-and-add
expansions in place of \texttt{mulw} and \texttt{divu}.

\section{Per-module comparison}

Table~\ref{tab:per-module-cmp} aligns the per-module counts using
the mapping in Table~\ref{tab:mapping}.

\begin{table}[ht]
\centering
\caption{Per-module instruction counts.  The reference's
  \texttt{xmss\_commons.o} is compared against the combined total of
  our \texttt{ltree}, \texttt{treehash}, and \texttt{bds.o} (where
  that combination makes sense).}
\label{tab:per-module-cmp}
\begin{tabular}{@{}lrrr@{}}
\toprule
Role & \texttt{impl/c/} & Reference & Delta \\
\midrule
SHA-2 core                   & 1953 & ---  & (delegated) \\
SHAKE / Keccak               & 1337 & 1259 & $-$78 \\
Hash wrappers                & 1082 &  684 & $-$398 \\
ADRS manipulation            &   99 &   29 & $-$70 \\
WOTS+                        &  760 &  498 & $-$262 \\
L-tree + treehash + BDS runtime & 1743 & $\subset$\,3169 & --- \\
BDS serialisation            &  473 &  --- & (not present) \\
XMSS / XMSS-MT core + commons & 1782 & 3682 & $+$1900 \\
Top-level API                & (in core) & 362 & --- \\
Parameters                   &  286 & 1077 & $+$791 \\
Utilities                    &   47 &   29 & $-$18 \\
\midrule
Total (archive-local)        & 9562 & 7620 & $-$1942 \\
\bottomrule
\end{tabular}
\end{table}

Several module-level asymmetries stand out.

\subsection{Parameters: 286 vs 1077}

Our \texttt{params.o} is roughly a quarter the size of the reference's.
Both modules do the same job---map a wire-format OID to the numeric
parameters $(n, w, h, d, \ldots)$---but the reference uses a large
\texttt{switch}/\texttt{if} cascade over all twelve standardised OIDs,
emitting many compare-and-branch sequences, whereas our
\texttt{params.c} uses a compact lookup table indexed by OID.  The
difference is entirely structural and is well understood as a C-level
coding choice: a jump table is cheaper on RV64 than a linear search.

\subsection{Hash wrappers and ADRS: the cost of explicit serialisation}

Our \texttt{xmss\_hash.o} (1082) and \texttt{address.o} (99) together
total 1181 instructions against the reference's \texttt{hash.o}
(684)~+~\texttt{hash\_address.o} (29)~=~713.

Our hash wrappers are deliberately more elaborate.  The reference
passes ADRS as \texttt{uint32\_t addr[8]} and serialises it to bytes
inside each hash wrapper via a shared \texttt{addr\_to\_bytes} helper.
We pre-serialise ADRS into a stack buffer and pass a pointer, which
(a) makes the hash wrappers match the shape of the Jasmin port we
intend to write, and (b) gives the compiler a stable input layout so
the PRF/F/H cores can be emitted with fixed operand sizes.  The cost
is visible: ~400 more instructions in the hash wrappers.  This is a
deliberate trade that pays off in the Jasmin port rather than in this
archive's instruction count.

\subsection{WOTS+: 760 vs 498}

Our \texttt{wots.o} is 262 instructions larger.  A portion of this
comes from the address-by-pointer calling convention (every F call
takes a pre-serialised ADRS; the reference passes \texttt{uint32\_t[8]}
and serialises inside F).  The rest is that our WOTS+ chain routine
explicitly checks chain-length invariants (\texttt{start + steps <= W-1})
and has a dedicated fast path for \texttt{steps == 0} that the
reference lacks.  These defensive checks are motivated by the Jasmin
port, where out-of-range chain indices would manifest as subtle
verification-side bugs.

\subsection{XMSS core and BDS: the shape differs, the size matches}

The reference bundles WOTS+ pk generation, L-tree, treehash, BDS
state, auth-path computation, XMSS keygen/sign/verify, and XMSS-MT
all into \texttt{xmss\_core\_fast.o} (3169~insns).  We split these
into separate compilation units: \texttt{bds.o} (1130),
\texttt{treehash.o} (484), \texttt{ltree.o} (129), \texttt{xmss.o}
(637), \texttt{xmss\_mt.o} (1145), and \texttt{xmss\_commons} logic
inlined into callers.  The combined total, 3525, is within 11\% of
the reference's single file, and the difference is dominated by our
additional \texttt{bds\_serialize.o} (473 insns) for wire-format
serialisation of BDS state---a feature the reference's
\texttt{xmss\_core\_fast.c} lacks entirely.

\section{Zbb: what changes when bit-manip is enabled}

Both implementations were re-profiled with \texttt{-march=rv64gc\_zbb}.
The extent of Zbb uptake differs sharply, and the mix reveals the
structural SHA-2 asymmetry from \S\ref{sec:sha2-asym}.

\begin{table}[ht]
\centering
\caption{Zbb uptake with \texttt{-march=rv64gc\_zbb}.}
\label{tab:zbb-cmp}
\begin{tabular}{@{}lrrr@{}}
\toprule
Metric                         & \texttt{impl/c/} & Reference & Delta \\
\midrule
I   (rv64gc)                   & 9505 & 7541 & \\
I   (rv64gc\_zbb)              & 9281 & 6982 & \\
\quad I removed                & $-$224 & $-$559 & \\
Zbb (rv64gc\_zbb)              &  164 &  132 & \\
Total (rv64gc\_zbb)            & 9502 & 7193 & \\
\quad Net instruction saving   & $-$60 & $-$427 & \\
\bottomrule
\end{tabular}
\end{table}

The reference achieves a 7$\times$ larger net code-size reduction from
Zbb than we do (427 vs 60 instructions), despite emitting 32 fewer Zbb
instructions than we do.  That inversion is explained by the mix:

\begin{table}[ht]
\centering
\caption{Zbb mnemonic mix with \texttt{-march=rv64gc\_zbb}.}
\label{tab:zbb-mix}
\begin{tabular}{@{}lrrll@{}}
\toprule
Mnemonic        & \texttt{impl/c/} & Reference & Native RV64 cost it replaces & Where \\
\midrule
\texttt{rolw}   & 41 & 0  & 3~insns (\texttt{srliw}+\texttt{slliw}+\texttt{or})  & 32-bit SHA-2 rotation \\
\texttt{roriw}  & 17 & 0  & 3~insns (same)                                      & 32-bit SHA-2 rotation \\
\texttt{rev8}   & 19 & 0  & multi-insn byte-swap                                & SHA-2 endianness \\
\texttt{rol}    & 17 & 38 & 3~insns (\texttt{srl}+\texttt{sll}+\texttt{or})     & 64-bit Keccak rho \\
\texttt{rori}   & 27 & 20 & 3~insns (same)                                      & 64-bit Keccak rho \\
\texttt{andn}   & 28 & 51 & 2~insns (\texttt{not}+\texttt{and})                 & Keccak chi, XMSS-MT masks \\
\texttt{maxu}   & 12 & 20 & $\sim$3~insns (branch-based)                        & BDS height comparisons \\
\texttt{minu}   &  3 &  2 & $\sim$3~insns (branch-based)                        & BDS/SHAKE min \\
\texttt{zext.h} &  0 &  1 & 2~insns (\texttt{slli}+\texttt{srli})               & 16-bit widening \\
\midrule
Total           & 164 & 132 & \multicolumn{2}{c}{} \\
\bottomrule
\end{tabular}
\end{table}

Two observations frame the comparison:

\begin{enumerate}[nosep]
  \item \textbf{Our archive has 32-bit SHA-2 rotations to optimise
    (\texttt{rolw}, \texttt{roriw}, \texttt{rev8});} the reference
    does not, because the reference's SHA-2 lives in OpenSSL's
    \texttt{libcrypto} and is not part of the analysed archive.  The
    77 Zbb instructions emitted in our \texttt{sha2\_local.c.o}
    (41~\texttt{rolw}~+~17~\texttt{roriw}~+~19~\texttt{rev8}) have no
    reference counterpart in this analysis: whatever Zbb
    \texttt{libcrypto} does or does not use is invisible here.
  \item \textbf{The larger net saving on the reference side is an
    archive-size effect, not a Keccak-shape effect.}  An earlier draft
    of this section attributed the reference's larger net saving to
    its Keccak rho step folding into \texttt{rol}/\texttt{rori} more
    idiomatically than ours.  Direct measurement (hand-unrolling our
    \texttt{keccak\_f1600} into the same straight-line form as the
    reference's \texttt{fips202.c}, then reprofiling) disproved that:
    the permutation itself is comparably sized in both, and the
    unrolled form actually inflated our archive overall.  The real
    source of the per-instruction asymmetry is elsewhere in the SHAKE
    module: we ship an incremental \texttt{\_ctx\_\{init,absorb,finalize,squeeze\}}
    API (driven by Jasmin-portability rules that forbid VLAs and
    discourage caller-allocated whole-message buffers), which the
    reference does not.  The extra API surface raises our baseline;
    the reference's smaller baseline makes any given Zbb saving look
    bigger in relative terms.  See \texttt{reports/jasmin-scope-hash-vs-upper.md}
    for the full design-tax accounting.
\end{enumerate}

The key qualitative finding is shared: \textbf{both} implementations
concentrate their Zbb uptake in the hash layer.  In our archive
142/164~=~87\% of Zbb lands in \texttt{sha2\_local}~+~\texttt{shake\_local};
in the reference, 108/132~=~82\% lands in \texttt{fips202.o}.  The
algorithm layer, in both cases, sees only opportunistic
\texttt{maxu}/\texttt{minu}/\texttt{andn} uses that do not affect
correctness or performance materially.

\section{Takeaways}

\begin{description}[style=nextline,leftmargin=1.5cm]
  \item[Same ISA surface.]
    Both implementations, compiled with \texttt{rv64gc}, use only
    I~+~M.  XMSS at the algorithm level is I-dominated and does not
    structurally require anything more.  M usage is below 2\% in both
    cases and comes entirely from compiler-emitted index/offset
    arithmetic.

  \item[Netting out SHA-2 equalises the archives.]
    The 1942-instruction gap between the two archives is almost
    exactly our bundled SHA-2 (\texttt{sha2\_local.c.o},
    1953~insns).  Everything else---hash wrappers, ADRS handling,
    WOTS+, BDS, XMSS---differs by at most a few hundred instructions
    per module.

  \item[Structural differences are design choices, not algorithm
    cost.]
    Our \texttt{params.o} is 791~insns smaller because we use a
    table; our hash and WOTS modules are several hundred insns
    larger because we pre-serialise ADRS and include defensive
    range checks.  These differences reflect the Jasmin-port
    preparation of our implementation, not anything about the XMSS
    algorithm itself.

  \item[Zbb matters most where SHA-2 lives.]
    Our archive benefits from Zbb via 32-bit SHA-2 rotations and
    \texttt{rev8}; the reference benefits via 64-bit Keccak
    rotations.  In both cases $\geq 82\%$ of Zbb uptake is in the
    hash layer.  The algorithm layer's Zbb benefit is negligible.
    For a self-contained XMSS implementation (ours), Zbb trims
    about 60 instructions; for the reference-as-archived (SHA-2
    opaque), it trims about 427 largely by improving Keccak.  A
    reference build that included SHA-2 in source form would show
    a far larger absolute Zbb benefit, dominated by the SHA-2
    compression core.

  \item[ISA guidance for a Jasmin RISC-V port is unchanged.]
    Both analyses point to \texttt{rv64im} as the minimum functional
    target and \texttt{rv64gc\_zbb} as the optimisation target for
    the hash layer.  Algorithm-layer Jasmin code can be written in
    portable RV64I~+~M without loss, and the hash layer is where
    ISA-specific work pays back.
\end{description}

\section{Reproducing this report}

This report is a synthesis of two independently generated reports.
To regenerate the raw data:

\begin{verbatim}
  # Our C implementation
  (cd impl/c && make rv)
  isa/scripts/analyse.sh       # → xmss_rv64_isa_profile.md
  isa/scripts/compare_zbb.sh   # → rv64gc vs rv64gc_zbb comparison

  # Upstream reference implementation
  isa/scripts/build_reference_lib.sh rv64gc
  isa/scripts/build_reference_lib.sh rv64gc_zbb
  isa/scripts/analyse.sh isa/binaries/libxmss-ref-rv64gc.a \
                         isa/reports/xmss_rv64_isa_profile_reference.md
  isa/scripts/analyse.sh isa/binaries/libxmss-ref-rv64gc_zbb.a \
                         /tmp/xmss_rv64_isa_profile_reference_zbb.md
\end{verbatim}

The tables in this report transcribe numbers from those generated
\texttt{.md} files and from the companion reports
\texttt{xmss\_rv64\_isa\_profile.tex} and
\texttt{xmss\_rv64\_isa\_profile\_reference.tex}.

\end{document}
